% Original analytic teaching diagram, 2026-11-09.
% 割线条件在 n=2 时确定的部分与 n=4 时待定的元素：s=(24,10)，y=(12,18)，单位 mm。
\definecolor{VB}{HTML}{7188A3}\definecolor{VBF}{HTML}{E5EAF0}
\definecolor{VG}{HTML}{829C8B}\definecolor{VGF}{HTML}{E8EEE9}
\definecolor{VY}{HTML}{B4A16F}\definecolor{VYF}{HTML}{F1ECDD}
\definecolor{VR}{HTML}{A7787D}\definecolor{VRF}{HTML}{F0E5E6}
\definecolor{VP}{HTML}{9588A5}\definecolor{VPF}{HTML}{ECE8F0}
\definecolor{VI}{HTML}{3B4252}\definecolor{VM}{HTML}{6B7280}\definecolor{VL}{HTML}{4C4D4F}
\begin{tikzpicture}[x=1mm,y=1mm,>=Stealth,font=\figurefont\fontsize{8}{10}\selectfont,
lab/.style={text=VM,font=\figurefont\fontsize{7.5}{9}\selectfont},
vec/.style={->,line width=1.1pt},cell/.style={minimum size=6.4mm,inner sep=0pt,draw=VL,line width=.5pt}]
\begin{scope}
\draw[vec,VG] (0,0)--(24,10) node[right,text=VI] {$\bm{s}$};
\draw[vec,VY] (0,0)--(-7.7,18.5) node[above,text=VI] {$\bm{u}$};
\fill[VL] (0,0) circle (.5);
\node[text=VI] at (10,-7) {输入};
\end{scope}
\draw[->,VL,line width=.75pt] (29,14)--(41,14) node[midway,above,text=VI] {$\bm{B}$};
\begin{scope}[shift={(56,0)}]
\draw[VR,line width=.95pt,dash pattern=on 3pt off 1.8pt] (-2.5,30)--(10.8,-1.9);
\draw[vec,VB] (0,0)--(12,18) node[right,text=VI] {$\bm{y}$};
\fill[VL] (0,0) circle (.5);
\node[text=VI,anchor=west] at (-1,30.5) {$\bm{Bu}$可落的直线};
\node[text=VI] at (8,-7) {输出};
\end{scope}
\node[text=VI] at (38,-14) {(a)};
\begin{scope}[shift={(77,1.5)}]
\foreach \i in {0,1,2,3}{\foreach \j in {0,1,2,3}{
 \pgfmathtruncatemacro{\k}{(\i==0 || \j==0) ? 1 : 0}
 \ifnum\k=1 \node[cell,fill=VBF] at (6.4*\j,6.4*3-6.4*\i) {}; \else \node[cell,fill=VGF] at (6.4*\j,6.4*3-6.4*\i) {}; \fi
}}
\node[text=VI] at (9.6,-15.5) {(b)};
\end{scope}
\end{tikzpicture}
